\section{Introduction}

Computability theory became possible once
precise models became available for modelling the
commonplace phenomenon of mechanical calculation.
The theory that evolved has been used to explain
human experience and to suggest how artificial
computing devices should be built.  It is also
worth studying for its own sake.

The commonplace phenomenon of learning surely
merits similar attention.  The problem is to discover good models that are interesting to study
for their own sake and promise to be relevant both
to explaining human experience and to building
devices that can learn.  The models should also
shed light on the limits of what can be learnt,
just as computability does on what can be computed.

In this paper we shall say that a program for
performing a task has been acquired by \emph{learning} if
it has been acquired by any means other than
explicit programming.  Among human skills some
clearly appear to have a genetically preprogrammed
element while some others consist of executing an
explicit sequence of instructions that has been
memorized.  There remains a large area of skill
acquisition where no such explicit programming is
identifiable.  It is this area that we describe
here as learning.  The recognition of familiar
objects, such as tables, provide such examples.
These skills often have the additional property
that, although we have learnt them, we find it
difficult to articulate what algorithm we are
really using.  In these cases it would be especially significant if machines could be made to
acquire them by learning.

This paper is concerned with precise computational models of the learning phenomenon.  We shall
restrict ourselves to skills that consist of recognizing whether a \emph{concept} (or predicate) is true
or not for given data.  We shall say that a concept Q has been learnt if a program for recognizing it has been deduced (i.e.\ by some method
other than the acquisition from the outside of the
explicit program).

The main contribution of this paper is that it
shows that it is possible to design \emph{learning
machines} that have all three of the following properties.

\begin{enumerate}
\item[(A)] The machines can provably learn whole
classes of concepts.  Furthermore these classes can
be characterized.

\item[(B)] The classes of concepts are appropriate
and nontrivial for general purpose knowledge, and

\item[(C)] The computational process by which the
machines deduce the desired programs requires a
feasible (i.e.\ polynomial) number of steps.
\end{enumerate}

A learning machine consists of a \emph{learning protocol} together with a \emph{deduction procedure}.  The
former specifies the manner in which information is
obtained from the outside.  The latter is the
mechanism by which a correct recognition algorithm
for the concept to be learnt is deduced.  At the
broadest level the suggested \emph{methodology} for
studying learning is the following:  define a
plausible learning protocol and investigate the
class of concepts for which recognition programs
can be deduced in polynomial time using the
protocol.

The specific protocols considered in this
paper allow for two kinds of information supply.
First the learner has access to a supply of typical data that positively exemplify the concept.
More precisely, it is assumed that these positive
examples have a probabilistic distribution determined arbitrarily by nature.  A call of a routine
EXAMPLES produces one such positive example.  The
relative probability of producing different examples is determined by the distribution.  The
second source of information that may be available is a routine ORACLE.  In its most basic
version, when presented with data it will tell the
learner whether or not the data positively exemplifies the concept.

The major remaining design choice that has to
be made is that of knowledge representation.
Since our declared aim is to represent general
knowledge, it seems almost unavoidable that we use
some kind of logic rather than, for example,
formal grammars or geometrical constructs.  In
this paper we shall represent concepts as Boolean
functions of a set of propositional variables.
The recognition algorithms that we attempt to
deduce will be therefore Boolean circuits or
expressions.

The adequacy of the propositional calculus
for representing knowledge in practical learning
systems is clearly a crucial question.  Few would
argue that this much power is not necessary.  The
question is whether it is enough.  There are
several arguments suggesting that such a system
would, at least, be a very good start.  First,
when one examines the most famous examples of
systems that embody preprogrammed knowledge,
namely expert systems such as DENDRAL and MYCIN,
essentially no logical notation beyond the propositional calculus is used.  Surely it would be
over-ambitious to try to base learning systems
on representations that are more powerful than
those that have been successfully managed in programmed systems.  Second, the results in this
paper can be negatively interpreted to suggest
that the class of learnable concepts even within
the propositional calculus is severely circumscribed.  This suggests that the search for extensions to the propositional calculus that have
substantially larger learnable classes may be a
difficult one.

The positive conclusions of this paper are
that there are specific classes of concepts that
are learnable in polynomial time using learning
protocols of the kind described.  These classes
can all be characterized by defining the class
of programs that recognize them.  In each case
the programs are special kinds of Boolean expressions.  The three classes are:  (i) conjunctive normal form expressions with a bounded
number of literals in each clause, (ii) monotone
disjunctive normal form expressions, and (iii)
arbitrary expressions in which each variable
occurs just once.  In the first of these no calls
of the oracle are necessary.  In the last no
access to typical examples is necessary but the
oracles need to be more powerful than the one
described above.

The deduction procedure will in each case output an expression that with high likelihood closely
approximates the expression to be learnt.  Such an
approximate expression never says yes when it
should not, but may say no on a small fraction of
the probability space of positive examples.  This
fraction can be made arbitrarily small by increasing the runtime of the deduction procedure.  Perhaps the main technical discovery contained in the
paper is that with this probabilistic notion of
learning highly convergent learning is possible for
whole classes of Boolean functions.  This appears
to distinguish this approach from more traditional
ones where learning is seen as a process of ``inducing'' some general rule from information that is
insufficient for a reliable deduction to be made.
The power of this probabilistic viewpoint is
illustrated, for example, by the fact that an
arbitrary set of polynomially many positive examples cannot be relied on to determine expressions
consisting of even just a single monomial in any
reliable way.

There is another aspect of our formulation that
is worth emphasizing.  In a learning system that is
about to learn a new concept there may be an
enormous number of propositional variables available.  These may be primitive inputs, the values of
preprogrammed concepts, or the values of concepts
that have been learnt previously.  We want the
complexity of learning the new concept to be related only to the number of variables that may be set
in natural examples of it, and not on the cardinality of the universe of available variables.
Hence the questions asked of ORACLE and the values
given by EXAMPLES will not be truth assignments to
all the variables.  In the archetypal case they
will specify an assignment to a subset of the
variables that is still sufficient to guarantee
the truth of the function.

Whether the classes of learnable Boolean concepts can be extended significantly beyond the
three classes given is an interesting question.
There is circumstantial evidence from cryptography,
however, that the whole class of functions computable by polynomial size circuits is not
learnable.  Consider a cryptographic scheme that
encodes messages by evaluating the function $E_k$
where $k$ specifies the key.  Suppose that this
scheme is immune to chosen plaintext attack in the
sense that even if the values of $E_k$ are known for
polynomially many different inputs, it is computationally infeasible to deduce an algorithm for $E_k$
or for an approximation to it.  This is equivalent
to saying, however, that $E_k$ is not learnable.
The conjectured existence of good cryptographic
functions that are easy to compute therefore implies that some easy to compute functions are not
learnable.

If the class of learnable concepts is as
severely limited as suggested by our results then
it would follow that the only way of teaching more
complicated concepts is to build them up from such
simple ones.  Thus a good teacher would have to
identify, name and sequence these intermediate
concepts in the manner of a programmer.  The results of \emph{learnability theory} would then indicate
the maximum granularity of the single concepts that
can be acquired without programming.

In summary, this paper attempts to explore the
limits of what is learnable as allowed by algorithmic complexity.  The results are distinguishable
from the diverse body of previous work on learning
because they attempt to reconcile the three properties (A)--(C) mentioned earlier.  Closest in
rigour to our approach is the inductive inference
literature (see Angluin and Smith \cite{1} for a survey)
that deals with inducing such things as recursive
functions or formal grammars (but not Boolean
functions) from examples.  There is a large body
of work on pattern recognition and classification,
using statistical and other tools, (e.g.\ \cite{4}) but
the question of general knowledge representation is
not addressed there directly.  Learning, in various
less formal senses, has been studied widely as a
branch of artificial intelligence.  A survey and
bibliography can be found in \cite{2,6}.  In their
terminology the subject of this paper is concept
learning.
